% =====================================================================
%  Methods chapter for an MDPI-style journal manuscript
%  Tool: QGIS plugin for 3D implicit geological modelling (GemPy-based)
%
%  Compiles standalone with pdflatex (article class). To use the MDPI
%  template instead, replace the preamble with \documentclass[...]{Definitions/mdpi}
%  and move the \section blocks into the manuscript body.
% =====================================================================
\documentclass[11pt,a4paper]{article}

\usepackage[utf8]{inputenc}
\usepackage[T1]{fontenc}
\usepackage{amsmath,amssymb}
\usepackage{bm}
\usepackage[margin=2.5cm]{geometry}
\usepackage{hyperref}

\title{Materials and Methods: A QGIS Plugin for Three-Dimensional
Implicit Geological Modelling and Uncertainty Estimation}
\author{}
\date{}

\begin{document}
\maketitle

% =====================================================================
\section{Materials and Methods}
% =====================================================================

\subsection{Software Architecture and Input Data}

The tool is implemented as a plugin for QGIS ($\geq$3.40) written in
Python 3, and couples the geographic-data handling of QGIS with the
open-source implicit geomodelling engine GemPy \cite{Gempy_2021}. It is
designed to let a user move directly from georeferenced field data held
in a QGIS project to a three-dimensional structural model and its
associated uncertainty, without leaving the GIS environment.

Three input layers are required and are read directly from the active
QGIS project: (i) a \emph{surface-point} layer, in which each vertex
carries planar coordinates $(X, Y)$, an elevation $Z$ and a categorical
\emph{formation} attribute; (ii) an \emph{orientation} layer, in which
each measurement carries a location, a dip and a dip-azimuth (or,
equivalently, the components of the unit pole vector) and a polarity; and
(iii) a digital terrain model (DTM) supplied as a raster. An optional
polygon layer may be used to define an area of interest (AOI) that clips
the modelling domain. Vector layers in any QGIS-supported format are
internally converted to a tabular representation, and all layers are
required to share a common coordinate reference system (CRS) so that the
geological observations and the topography are co-registered.

\subsection{Implicit Structural Modelling}

The structural model is computed with GemPy, which reconstructs
geological interfaces as iso-surfaces of a continuous scalar (potential)
field interpolated from the surface points and orientation data by a
co-kriging scheme \cite{Gempy_2021}. Formations are organised into ordered
\emph{structural groups} whose stratigraphic and tectonic relationships
are specified by the user as \texttt{ERODE}, \texttt{ONLAP},
\texttt{FAULT} or \texttt{BASEMENT} relations; faults are detected
automatically from formation names and may offset the remaining series
according to configurable fault relations. The modelling domain is
discretised on a regular grid whose horizontal extent is derived from the
DTM (or the AOI) and whose vertical extent spans from a user-defined basal
depth to the maximum terrain elevation. The DTM is incorporated as the
upper boundary of the model, so that geological surfaces are truncated at
the topography. After interpolation, the explicit triangulated surfaces of
each interface are extracted from the scalar field and rendered
interactively in an embedded PyVista \cite{Sullivan2019} viewer.

\subsection{Mesh Generation and Export}

Each geological interface and the topographic surface can be exported as a
watertight triangular mesh in STL format for downstream use (e.g., 3D
printing or numerical simulation). The topographic solid is generated by
Delaunay triangulation of the DTM nodes and vertical extrusion to the
model base, while the structural surfaces are taken from the interpolated
iso-surfaces. Prior to export, all meshes are passed through a cleaning
and repair stage, coincident-vertex merging, consistent normal
orientation, and topological repair with the MeshFix
algorithm \cite{Attene2010MeshFix}, and the terrain is optionally smoothed
(Gaussian filtering of the elevation field only, preserving the
planimetric footprint and hence the CRS). This step removes non-manifold
edges, self-intersections and folded faces that would otherwise compromise
meshing downstream.

\subsection{Uncertainty Estimation by Monte Carlo Perturbation}

Because orientation measurements are the input most affected by
observational error, the tool quantifies structural uncertainty through a
Monte Carlo ensemble in which the orientation data are perturbed and the
model is recomputed \cite{wellmann}. For each of $N$ realisations,
every orientation is independently perturbed by adding Gaussian noise to
its dip and dip-azimuth,
%
\begin{equation}
\delta_k \sim \mathcal{N}\!\left(\delta_0,\sigma_\delta^2\right),
\qquad
\alpha_k \sim \mathcal{N}\!\left(\alpha_0,\sigma_\alpha^2\right),
\end{equation}
%
where $\delta_0$ and $\alpha_0$ are the measured dip and azimuth, and
$\sigma_\delta$ and $\sigma_\alpha$ are user-defined standard deviations
(default $5^\circ$ and $10^\circ$, respectively); dips are
clipped to the physically admissible range and azimuths are wrapped modulo
$360^\circ$. The full model is rebuilt and recomputed for each
realisation, keeping the geometry, grid and structural configuration
fixed, so that the only source of variability is the orientation noise.

Each realisation $k$ yields a discrete lithology label $l_i^{(k)}$ for
every grid cell $i$. Aggregating across the $N$ realisations gives, for
every cell, the relative frequency of each lithology $l$,
%
\begin{equation}
p_i(l) = \frac{1}{N}\sum_{k=1}^{N}
\mathbb{1}\!\left[\,l_i^{(k)} = l\,\right].
\end{equation}
%
From this distribution the tool derives, per cell: the \emph{modal
lithology} $\hat{l}_i = \arg\max_l p_i(l)$; the \emph{modal probability}
%
\begin{equation}
P_i = \max_l p_i(l),
\end{equation}
%
a confidence measure that approaches~1 where realisations agree and
decreases where they diverge; and the \emph{information (Shannon) entropy}
%
\begin{equation}
H_i = -\sum_{l} p_i(l)\,\ln p_i(l),
\end{equation}
%
which is zero where a single lithology is always predicted and maximal
where several lithologies are equally likely~\cite{ref-wellmann}. The
resulting probability, entropy and modal-lithology volumes are written to
disk as gridded arrays (together with the model extent) and can be
explored as interactive cross-sections. Summary statistics---the mean
modal probability, the mean entropy, and the fraction of cells whose modal
probability falls below a confidence threshold ($0.8$)---are reported to
characterise the overall reliability of the model.

% =====================================================================
\begin{thebibliography}{9}

\bibitem{ref-gempy}
de la Varga, M.; Schaaf, A.; Wellmann, F.
GemPy 1.0: open-source stochastic geological modeling and inversion.
\emph{Geosci. Model Dev.} \textbf{2019}, \emph{12}, 1--32.

\bibitem{ref-wellmann}
Wellmann, J.F.; Regenauer-Lieb, K.
Uncertainties have a meaning: Information entropy as a quality measure for
3-D geological models.
\emph{Tectonophysics} \textbf{2012}, \emph{526--529}, 207--216.

\bibitem{ref-pyvista}
Sullivan, C.B.; Kaszynski, A.
PyVista: 3D plotting and mesh analysis through a streamlined interface for
the Visualization Toolkit (VTK).
\emph{J. Open Source Softw.} \textbf{2019}, \emph{4}, 1450.

\bibitem{ref-meshfix}
Attene, M.
A lightweight approach to repairing digitized polygon meshes.
\emph{Vis. Comput.} \textbf{2010}, \emph{26}, 1393--1406.

\end{thebibliography}

\end{document}
